% Imporditud: tools/import_eksperimendid.py
% Allikas: Eesti füüsikaolümpiaadi eksperimendi ülesannete
%   kogu (Rael Kalda, 2023), ülesanne Ü176, lahendus L176.
\setAuthor{Jaan Kalda/Eero Uustalu}
\setRound{lõppvoor}
\setYear{2018}
\setNumber{G E2}
\setDifficulty{5}
\setTopic{elektriahelad}
\setVahendid{kolme väljundklemmiga must kast, mille sees on takisti, võrdlemisi suure sisetakistusega patarei ja kondensaator; multimeeter; stopper; juhtmed; millimeeterpaber}
\setPunktid{14}
\setAllikas{Eksperimendi kogu Ü176, lahendus lk 135}
\setLahendusseis{toores}

\prob{Must kast}
Tehke kindlaks mustas kastis olev elektriline skeem ja määrake takisti takistus, patarei sisetakistus ning kondensaatori mahtuvus.

\hint

\solu
Kolm elementi on mustas kastis tähtühenduses. Sellest arusaamiseks paneme tähele, et ainult ühe klemmipaari puhul ei lähe vool pika aja möödudes nulli. Märgime ära stabiliseeruva voolutugevuse I nende klemmide vahel. Kasutades multimeetrit voltmeetrina, mõõdame nende klemmide vahelise pinge - see võrdub patarei elektromotoorjõuga E. Nendest mõõtmistest saame leida Ohmi seaduse abil takisti takistuse ja patarei sisetakistuse summa:

E R+r =

I

Eraldi takistuse leiame, kui laadida esmalt kondensaator täiesti täis (ampermeet-

riga saab kontrollida laadimisvoolu vähenemist) ja siis ühendada läbi ampermeet-

ri takistile (st ühendada ampermeetriga kaks vastavat väljundklemmi): esimese

hetke lugem annab voolu, kui takistil on patarei pinge. Kondensaatori mahtuvuse

leidmiseks mõõdame stopperiga aega ning konstrueerime tühjenemisvoolu graa-

fiku I (t). Selle graafiku alune pindala võrdub kondensaatori algse laenguga Q. Siit

kondensaatori mahtuvus C = Q/U . Alternatiivselt saame kondensaatori mahtu-

vuse, kui mõõdame voolutugevuse kahel ajahetkel ja kasutame seost kondensaat tori tühjenemisvoolu kohta: I = $I_0$ e$-$ RC .
\probend
